Non-fatal (surviving) cases of aortic pathology associated with gaseous findings (HPVG / pneumatosis) and pain, plus notes on chronic aortic pain.
Documented non-lethal human cases linking aortic disease with portal/mesenteric gas and pain are clinically rare. Most published aortic + HPVG associations are acute and high-mortality. The following are the clearest surviving examples.
Confirmed non-fatal aortic + gas + pain cases
Stanford B aortic dissection + post-TEVAR HPVG (full recovery, long-term survival)
46-year-old man presented with acute chest pain and was diagnosed with Stanford type B aortic dissection. After thoracic endovascular aortic repair (TEVAR), he developed acute abdominal pain and distension. Abdominal CT showed hepatic portal venous gas (HPVG).
Managed entirely conservatively (gastrointestinal decompression, fluid resuscitation, electrolyte correction, antibiotics, anti-inflammatory agents, anticoagulation). Symptoms resolved; follow-up CT showed disappearance of HPVG. Discharged on day 10.
5-year follow-up: no recurrent physical discomfort related to HPVG or residual symptoms. This is the clearest documented non-fatal aortic-dissection-associated HPVG case with pain.
HPVG + pneumatosis in a patient with prior aortic dissection (conservative survival)
An 86-year-old woman with dementia and a history of aortic dissection presented with vomiting and melena. CT showed HPVG and intestinal pneumatosis. Laparotomy was considered but withheld due to poor general condition. Conservative treatment (rehydration + antibiotics) led to resolution of HPVG and clinical recovery.
Broader HPVG series containing survivors (some with vascular comorbidities)
In modern cohorts of HPVG (not exclusively aortic), 60 % of patients recover. Conservative management succeeds in a substantial fraction when peritonitis, shock, and clear transmural necrosis are absent. Abdominal pain is the dominant symptom (≈88 %). Resolution of portal gas on follow-up CT typically occurs within days (median ≈3.8 days in one series). These series include patients with atherosclerotic or prior aortic disease, but pure chronic aortic + gas linkage remains sparse.
Additional non-aortic but instructive surviving HPVG/pneumatosis cases (e.g., after pancreaticoduodenal artery aneurysm embolization or superior mesenteric artery syndrome) confirm that gas findings + pain can resolve with conservative care when ischemia is limited or reversible.
Chronic aortic pain syndromes
True chronic aortic pain syndromes (pain lasting weeks to months) are well-described but are not characteristically associated with gaseous findings (HPVG or pneumatosis):
Chronic type B aortic dissection or aneurysmal expansion can produce persistent or recurrent chest, back, or abdominal pain due to wall stretch, progressive enlargement, or intermittent malperfusion.
Pain is often dull, constant, or migratory rather than the classic acute “tearing” pain.
Imaging (CTA/MRA) shows chronic dissection flaps, thrombosed false lumen, or expanding aneurysm; gas is absent unless an acute superimposed ischemic or infectious event occurs.
Management is usually medical (blood-pressure control) with surveillance; intervention is reserved for complications (expansion, recurrent pain refractory to medical therapy, malperfusion, or rupture signs).
No large series or consistent case literature links chronic aortic pain syndromes with ongoing or recurrent portal/mesenteric gas. When gas appears, the presentation is almost always acute (or post-procedural) and signals a new ischemic or barrier-disruption event rather than a chronic gaseous process.
Summary
Surviving documented patients with aortic pathology + gaseous findings + pain exist but are an uncommon diagnosis; the best-documented example is the post-TEVAR Stanford B dissection case with complete recovery and multi-year asymptomatic follow-up. Chronic aortic pain is a recognized clinical entity driven by mechanical wall stress or expansion, yet it occurs without the gaseous radiographic markers that characterize the acute ischemic pathway.

Documented human patients with aortic pathology (dissection, aneurysm, or related stress/ischemia) presenting with pain and gaseous findings (primarily hepatic portal venous gas [HPVG], pneumatosis intestinalis, or related portomesenteric gas).
These cases illustrate clinically observed intersections of aortic disease with acute gaseous phenomena, typically in the setting of mesenteric malperfusion, bowel ischemia/necrosis, or post-procedural complications. Abdominal (or chest) pain is a near-universal presenting feature. They correspond to the acute, catastrophic pathway distinguished in the repository (physical gas emboli secondary to barrier failure or ischemia), not the chronic molecular priming cascade.
Key cited case reports and series
Catastrophic STS/SVS type B3 aortic dissection + HPVG + multi-organ malperfusion
64-year-old hypertensive male presented with severe, unremitting abdominal pain (8 hours). CTA showed type B3 dissection with celiac trunk and SMA straddled by the flap (static + dynamic obstruction), diffuse HPVG, pneumatosis intestinalis, and multi-organ ischemia/necrosis. Fatal outcome. Described as a “deadly triad” of aortic dissection, mesenteric malperfusion, and transmural bowel necrosis.
Thoracic (Stanford B) aortic dissection + post-TEVAR HPVG
46-year-old hypertensive male presented with acute chest pain; diagnosed with Stanford B dissection. After thoracic endovascular aortic repair (TEVAR), developed acute abdominal pain and distension. Abdominal CT demonstrated HPVG. Managed conservatively (decompression, fluids, antibiotics, etc.); recovered and remained asymptomatic at 5-year follow-up. First reported association of HPVG after TEVAR for aortic dissection (attributed to possible reperfusion injury).
Thrombosed abdominal aortic aneurysm + extensive HPVG
78-year-old female presented with altered mental status; progressed to extensive HPVG, pneumatosis intestinalis, and huge abdominal aortic aneurysm on CT. Fatal course despite intensive care (death within ~12–36 hours of imaging).
Severe abdominal aortic atherosclerosis + classic pneumatosis intestinalis + HPVG
86-year-old man with multiple comorbidities presented with abdominal distension, tenderness, and absent bowel sounds (in context of dehydration/hyperkalemia/mental-status change). CT showed HPVG and pneumatosis intestinalis coexisting with severe atherosclerotic calcification of the abdominal aorta (suggestive of mesenteric atherosclerotic disease).
Portal pneumatosis + aortic endoleak after recent aortic valve surgery
75-year-old man (recent Bentall–De Bono aortic valve replacement) presented with nonspecific abdominal pain and poor general condition, hemodynamic instability, and severe lactic acidosis. CTA showed aortic endoleak plus extensive portal pneumatosis (gas in superior mesenteric vein, splenoportal axis, portal vein) corresponding to massive intestinal ischemia. Fatal within hours.
Pneumatosis intestinalis with portal/mesenteric/renal gas + disruption of abdominal aortic perfusion
Young patient with chronic colonic pseudo-obstruction developed cardiac arrest; CT after ROSC showed massive colonic dilatation, pneumatosis coli, extensive gas in mesenteric veins/arteries, massive portal gas, gas in renal veins, and disruption of abdominal aortic perfusion. Fatal.
Additional supporting observations  
RadioGraphics review of portomesenteric vein gas includes illustrative CT of mesenteric ischemia secondary to abdominal aortic dissection with gas in mesenteric veins.  
Older literature documents massive visceral necrosis after aortic dissection with intraportal, mesenteric-vein, and bowel-wall gas on CT, presenting with acute abdominal pain.  
Series of HPVG cases frequently list abdominal pain (≈88 % in one cohort of 25 patients) as the dominant symptom; intestinal ischemia/necrosis (often linked to aortic/mesenteric vascular disease or atherosclerosis) is the leading etiology.
Clinical pattern across these human cases
Pain: Severe, unremitting abdominal pain (or chest pain radiating/migrating) is the dominant presenting complaint in essentially all acute presentations.  
Aortic findings: Dissection (type B or involving visceral arteries), aneurysm (thrombosed or with endoleak), or severe atherosclerosis causing mesenteric malperfusion.  
Gaseous findings: HPVG, pneumatosis intestinalis, or portomesenteric gas—radiologic markers of mucosal barrier breach and/or bowel ischemia/necrosis.  
Outcome: High mortality when HPVG coexists with transmural ischemia and aortic malperfusion; occasional recovery with prompt recognition and supportive care (especially post-procedural).
These documented human cases confirm that aortic wall stress/malperfusion can clinically coincide with acute gaseous syndromes and severe pain. They remain rare under clinical observation, typically occurring only after a critical threshold of ischemia or barrier failure has already been crossed—consistent with the permissive, threshold-dependent character of the modeled pathways.
